\subsection{带算子空间}

设 X 是一个拓扑空间，G 是一个拓扑群。如果下列条件得到满足，就称 G 连续作用在 X 上：

\begin{enumerate}
\def\labelenumi{\arabic{enumi})}
\item
  X 以 G 为算子群；换言之，X 赋有一个外合成法则 \((s, x) \to s.x\)，G 是它的算子集，并且 \(s.(t.x) = (st).x\) 与 e.x = x 对所有 \(s, t \in G\) 与所有 \(x \in X\) 成立。
\item
  映射 \((s, x) \to s.x\)（从 \(\mathbf{G} \times \mathbf{X}\) 到 \(\mathbf{X}\) 中的）是连续的。
\end{enumerate}

引理 1 若拓扑群 G 连续作用在拓扑空间 X 上，则对每个 \(s \in \mathbf{G}\)，映射 \(x \to s.x\) 是 X 到 X 上的一个同胚。

因为这个映射是一个连续双射，其逆 \(x \rightarrow s^{-1}.x\) 也是连续的。

我们回顾，对每个 \(x \in X\)，x 在 G 的元素 s 下的变换 s.x 组成的集 G.x 称为 x 的轨道（关于算子群 G 的），而使 s.x = x 的所有 \(s \in G\) 组成的集是 G 的一个子群，称为 x 的稳定子群。关系 \(\mathrm{R}(x, y) : ``y \text{ 属于轨道 } x''\) 是 X 上的一个等价关系，称为由 G 定义的等价关系；关于这个关系的诸等价类就是 X 的诸点的轨道。拓扑空间 X/R 称为 X 的轨道空间（关于 G 的），或 X 关于群 G 的商空间，记作 X/G；并称 X/G 的拓扑为 X 的拓扑经 G 的商。

引理 2 若拓扑群 G 连续作用在拓扑空间 X 上，则由 G 定义的等价关系 R 是开的。

因为关于 \(\mathbf{R}\) 取的 \(\mathbf{U}\)（\(\mathbf{X}\) 的一个开子集）的饱和集是集合 \(\bigcup_{s\in G}s.\mathbf{U},\)，而由引理 1，每个 \(s.\mathbf{U}\) 都是开的。

例。1) 设 H 是拓扑群 G 的一个子群。H 通过外法则 \((s, x) \rightarrow sx\) 连续作用在 G 上。H 也通过外法则 \((x, s) \rightarrow sxs^{-1}\) 连续作用在 G 上。

\begin{enumerate}
\def\labelenumi{\arabic{enumi})}
\setcounter{enumi}{1}
\item
  *若 K 是一个拓扑除环（§ 6 第 7 目），则乘法群 K* 通过外法则 \((s, x) \to sx\) 连续作用在 K 上。*
\item
  设 G 是一个拓扑群，X 是一个拓扑空间。则映射 \((s, x) \to x\)（从 \(G \times X\) 到 X 中的）是 X 上的一个外合成法则，且 G 关于这个法则连续作用在 X 上；这时称 G 平凡地作用在 X 上。
\end{enumerate}

注。与其说拓扑群 G 连续作用在拓扑空间 X 上，人们常说 G 在 X 的左边连续作用。当 G 的反群 \(G^{0}\)（拓扑群）连续作用在 X 上时，我们说 G 在 X 的右边连续作用。这等同于说：X 有一个连续的外合成法则 \((s, x) \to s.x\)，以 G 为算子集，满足 \(s.(t.x) = (ts).x\) 与 e.x = x。这样的法则常写在右边：\((s, x) \to x.s\)（术语即由此而来），这时有 \((x.t).s = x.(ts)\)。若 G 通过法则 \((s, x) \to x.s\) 在 X 上右边连续作用，则 G 也依照外法则 \((s, x) \to x.s^{-1}\) 在 X 上左边连续作用，这是由公理 \((\mathrm{GT}_{\mathrm{II}})\) 得出的。

设 X（相应地，X'）是带算子群 G（相应地，G'）的集合，\(f: \mathbf{G} \to \mathbf{G}'\) 是一个同态，\(g: \mathbf{X} \to \mathbf{X}'\) 是一个映射。则称 \(f\) 与 \(g\) 是相容的，如果 \(g(s.x) = f(s).g(x)\) 对所有 \(s \in \mathbf{G}\) 与所有 \(x \in \mathbf{X}\) 成立。若 \(\mathbf{X}''\) 是第三个带算子群 G'' 的集合，\(f': \mathbf{G}' \to \mathbf{G}''\) 是一个同态，\(g': \mathbf{X}' \to \mathbf{X}''\) 是一个映射，且 \(f'\) 与 \(g'\) 相容，则 \(f' \circ f\) 与 \(g' \circ g\) 相容。当 X、X' 是拓扑空间且 G、G' 是分别连续作用在 X、X' 上的拓扑群时，就称 \((f, g)\) 为带算子空间 X 到带算子空间 X' 中的态射，只要 \(f\) 与 \(g\) 是连续且相容的。过渡到商，\(g\) 便诱导出连续映射 \(\mathbf{X}/\mathbf{G} \to \mathbf{X}'/\mathbf{G}'\)（第一章 § 3 第 4 目命题 6 的推论）。

设 G 是连续作用在拓扑空间 X 上的一个拓扑群，\(\varphi\) 是 X 到轨道空间 X/G 上的典范映射。设 A 是 X 的任一子集，\(A'\) 是 X 中作为 A 关于由 G 定义的等价关系 R 的饱和集的子空间（于是 \(A'\) 是 A 的诸点的轨道之并，称为 A 关于 G 的饱和集）。G 连续作用在 \(A'\) 上，其作用法则是 \((s, x) \to s.x\) 到 \(G \times A'\) 上的限制。此外，由于 R 是开的（引理 2）且 \(A'\) 是饱和的，由第一章 § 5 第 2 目命题 4 以及关系式 \(\varphi(A) = \varphi(A')\) 可得：

命题 10 子空间 \(\varphi(A)\)（\(X/G\) 的）到轨道空间 \(A'/G\) 上的典范双射是一个同胚。

现设 S 是 X 上的一个等价关系，使得对每个 \(s \in G\)，映射 \(x \to s.x\) 与 S 相容{[}换言之，使得关系 \(x \equiv y \pmod{S}\) 蕴含 \(s.x \equiv s.y \pmod{S}\){]}；为简便起见，我们将说关系 S 与群 G 相容。若 \(\psi\) 是 X 到 X/S 上的典范映射，且以 \(s.\psi(x)\) 表示 s.x 的模 S 类，则 X/S 关于外法则 \((s, \psi(x)) \to s.\psi(x) = \psi(s.x)\) 以 G 为算子群。此外：

命题 11 若 X 上的等价关系 S 是开的且与 G 相容，则 G 连续作用在 X/S 上。

由于 G 上的相等关系与 X 上的关系 S 都是开的，只需证明映射 \((s,x)\to s.\psi(x)=\psi(s.x)\)（从 \(G\times X\) 到 X/S 中的）是连续的（第一章 § 5 第 3 目命题 8 的推论）；而这由 \(\psi\) 的连续性以及映射

\[
(s, x) \rightarrow s. x.
\]

注。设 \(G'\) 是另一个连续作用在 X 上的拓扑群，并设 \(s.(s'x) = s'(s.x)\) 对所有 \(s \in G\)、\(s' \in G'\) 与 \(x \in X\) 成立。则等价关系 \(S\)（由 \(G'\) 定义的）与 \(G\) 相容，且由于 \(S\) 是开的（引理 2），可知 \(G\) 连续作用在 \(X/G'\) 上。类似地，\(G'\) 连续作用在 \(X/G\) 上。在这些情形下，我们说两个群 \(G\) 与 \(G'\) 在 \(X\) 上的两个作用可交换。

